% Copy-ready mathematical module. Requires amsmath, amssymb, amsthm,
% hyperref and theorem/lemma/proposition/corollary environments.
% No repository files are modified by this intake artifact.
\section{The complete first-half selector and its cyclotomic relations}
\label{sec:bp-increment}

The source dialogue is JT and ChatGPT, \emph{Boolean product formulation},
9--11 September 2026, exported 12 September 2026
(\href{https://chatgpt.com/share/6aa445c7-f990-83ea-a2b9-0a9cae13dc52}{shared source dialogue}).
The preserved transcript has SHA-256
\texttt{19d54ae24d712598}\allowbreak
\texttt{c14f222ee81eb19f1}\allowbreak
\texttt{92c9e944fbbfe500}\allowbreak
\texttt{9f91780709792f3}.
The locators below refer to the one-based lines of that export.
JT supplies the Boolean-product, cyclotomic, generalized-congruence and
adversarial directions; the following arguments state the mathematical
claims with their exact scopes. They make no blanket claim of novelty.
In particular no universal Erd\H{o}s--Straus theorem or counterexample is
asserted. The first-half proof and the explicit complete permutation fibres
below fill out the relation between the retained full-range code and the
dialogue's first-half code.

Write \(\operatorname{ES}(p)\) for the existence of positive integers
\(x,y,z\) such that \(4/p=1/x+1/y+1/z\).
Prime factorization, quadratic reciprocity with its supplementary laws,
and cyclicity of the multiplicative group of a finite field are used in
their usual classical forms. Every additional arithmetic implication in
this section is proved here.

\subsection{The denominator bound, divisor maps and exact code}

\begin{theorem}[First-half denominator bound]\label{thm:bp-first-half}
Let \(p\) be prime and \(p\equiv1\pmod4\). The minimum denominator in
every positive integer Erd\H{o}s--Straus solution lies strictly between
\(p/4\) and \(p/2\).
\end{theorem}
\begin{proof}
Retain a permutation of the original coordinates and name its ordered
values \(x\le y\le z\). Positivity gives \(1/x<4/p\), so \(x>p/4\).
Suppose \(x>p/2\); equality is impossible because \(p\) is odd.
Then
\[
 \frac4p\le\frac1x+\frac2y<\frac2p+\frac2y,
\]
so \(y<p\). Neither \(x\) nor \(y\) is divisible by \(p\).
The integer equation \(4xyz=p(xy+xz+yz)\) gives \(p\mid z\).
Write \(z=pk\). The original equality becomes
\[
 (4k-1)xy=pk(x+y).
\]
Thus \(t=(4k-1)/p\) is a positive integer,
\(t\equiv3\pmod4\), \(\gcd(t,k)=1\), and \(txy=k(x+y)\).
Write \(x=g\alpha,y=g\beta\), where \(g=\gcd(x,y)\) and
\(\gcd(\alpha,\beta)=1\). Then
\[
 tg\alpha\beta=k(\alpha+\beta),\qquad
 \gcd(\alpha\beta,\alpha+\beta)=1.
\]
Hence \(k=\alpha\beta D\) for an integer \(D>0\), and
\[
 tg=D(\alpha+\beta),\quad \gcd(t,D)=1,\quad
 pt=4\alpha\beta D-1.
\]
The inequality \(x>p/2\) implies
\[
 2D\alpha(\alpha+\beta)>4\alpha\beta D-1,
 \qquad 2D\alpha(\alpha-\beta)>-1.
\]
The last left side is an integer, so it is nonnegative and
\(\alpha\ge\beta\). Applying the same argument to \(y>p/2\)
gives \(\beta\ge\alpha\). Coprimality forces
\(\alpha=\beta=1\). Now \(tg=2D\) and \(\gcd(t,D)=1\)
give \(t\mid2\), contrary to \(t\equiv3\pmod4\).
The retained permutation recovers the original coordinate order.
\end{proof}

Fix henceforth a prime \(p=12h+1\), \(h\ge1\), and put
\[
 a=3h+j+1,\quad R=4a-p=4j+3,\quad S=pa,
 \qquad 0\le j<3h.
\]
Thus \(p/4<a<p/2\), \(3\le R\le p-2\), and
\(\gcd(R,pa)=1\). Indeed \(a<p\) is prime to \(p\), a common
divisor of \(a,R\) divides \(p\), and a common divisor of \(p,R\)
divides \(4a\); also \(R\) is odd. Define
\[
 E_a=\{u>0:u\mid a^2,\ R\mid4u+1\},\qquad
 M_a=\{u>0:u\mid a^2,\ R\mid4u+p\}.
\]
The middle condition is equivalently \(R\mid u+a\).

\begin{theorem}[Original ordered solutions and the bounded divisor code]
\label{thm:bp-code}
For \(u\mid a^2\), set
\[
 g=\gcd(a,u),\quad A=a/g,\quad B=u/g,\quad D=g^2/u.
\]
These are positive integers and define a bijection
\[
 \operatorname{Div}(a^2)\longleftrightarrow
 \{(A,B,D)>0:a=ABD,\ \gcd(A,B)=1\},\qquad u=B^2D.
\]
For \(\varepsilon\in\{0,1\}\), its hit condition is exactly
\[
 R\mid4u+p^\varepsilon
 \quad\Longleftrightarrow\quad R\mid A+p^{1-\varepsilon}B.
\]
On that set, define
\[
 C=\frac{A+p^{1-\varepsilon}B}{R},\qquad
 (x,y,z)=(ABD,p^\varepsilon ACD,pBCD).
 \tag{BP1}
\]
Then \(C>0\) is integral, \(4ABCD=A+p^{1-\varepsilon}B+pC\),
and (BP1) is an original ordered solution.

Put
\[
 M_H=6h,\quad J_H=3h,\quad\Lambda_H=36h^2,
 \quad N_H=216h^3=(p-1)^3/8.
\]
The map
\[
 c_H=(A-1)\Lambda_H+(p-1)j+2(B-1)+\varepsilon
 \tag{BP2}
\]
is a bijection from the rectangle
\(1\le A,B\le M_H,\ 0\le j<J_H,\ \varepsilon\in\{0,1\}\)
to \(0\le c_H<N_H\). On its decoded coordinates set
\[
 \chi_{H,p}(c_H)=
 \left\lfloor\frac{\gcd(4AB,p+4j+3)}{4AB}\right\rfloor
 \left\lfloor\frac1{\gcd(A,B)}\right\rfloor
 \left\lfloor\frac{\gcd(4j+3,A+p^{1-\varepsilon}B)}{4j+3}\right\rfloor.
\]
Then
\[
 \operatorname{ES}(p) \Longleftrightarrow\
 \exists\,0\le c_H<N_H:\chi_{H,p}(c_H)=1.
 \tag{BP3}
\]
\end{theorem}
\begin{proof}
Write \(a=gA,u=gB\) with \(\gcd(A,B)=1\). The divisibility
\(u\mid a^2\) implies \(B\mid gA^2\), hence \(B\mid g\).
Therefore \(D=g/B\) is integral and \(a=ABD,u=B^2D\).
Conversely these equalities give \(u\mid a^2\) with quotient
\(A^2D\), and \(\gcd(a,u)=BD\). This proves both inverse
compositions. Also \(A,B\le a\le M_H\).
Using \(4ABD=p+R\), one has
\[
 A(4u+p^\varepsilon)
 \equiv pB+p^\varepsilon A
 =p^\varepsilon(A+p^{1-\varepsilon}B)\pmod R.
\]
Both \(A\) and \(p\) are units modulo \(R\), proving the
biconditional. On its passing locus, multiplying the definition of
\(C\) by \(R=4ABD-p\) gives the tuple equation. The sum of the
three reciprocals in (BP1), over denominator \(pABCD\), has numerator
\(pC+p^{1-\varepsilon}B+A=4ABCD\).

For the code inverse divide \(c_H\) by \(\Lambda_H\) to recover
\(A-1\) and remainder \(v\), divide \(v\) by \(p-1\) to recover
\(j\) and remainder \(r\), then divide \(r\) by 2 to recover
\(B-1\) and \(\varepsilon\). The ranges of successive remainders
are precisely the prescribed ranges. The gcd-floor factors equal
the divisibility or coprimality bits because their quotients are
positive and at most one. Thus a passing code gives (BP1).

For the converse, Theorem~\ref{thm:bp-first-half} selects a first
denominator \(a\) in this interval, retaining its original position.
For the other denominators define
\[
 d_y=Ry-S=Sy/z>0,\quad d_z=Rz-S>0,
 \qquad d_y d_z=S^2.
\]
Conversely a positive \(d_y\mid S^2\) with \(R\mid d_y+S\)
determines \(d_z=S^2/d_y\) and
\(y=(d_y+S)/R,z=(d_z+S)/R\); the second quotient is integral
because \(d_y\equiv-S\) is a unit modulo \(R\).
Since \(p\nmid a\), the exponent \(i=v_p(d_y)\) is 0,1 or 2.
For \(i=0\), set \(\varepsilon=0,u=a^2/d_y\); multiplying
\(d_y+pa\equiv0\) by \(4u/a\) gives \(p(1+4u)\equiv0\).
For \(i=2\), set \(\varepsilon=0,u=d_y/p^2\); multiplication
by \(4/p\) gives the same congruence, with the final two original
denominators exchanged in (BP1). For \(i=1\), write \(d_y=pv\)
and set \(\varepsilon=1,u=a^2/v\). Then \(v\equiv-a\), hence
\(u\equiv-a\), giving the middle condition. All three cases
therefore give a passing code. The exterior exchange is recorded,
not identified with the original ordering. This proves (BP3).
\end{proof}

\begin{proposition}[Radix embedding and the full permutation fibres]
\label{prop:bp-radix}
Retain the repository's full rectangle
\[
 J_F=6h,\quad M_F=9h,\quad\Lambda_F=108h^2,
 \quad N_F=972h^3,
\]
with the same shell formula and code
\(c_F=(A-1)\Lambda_F+18hj+2(B-1)+\varepsilon\).
The full-range hits satisfying \(a\le6h\) are bijective to the
first-half hits by decoding and recoding the same four coordinates.
The map preserves \(a,R,S,A,B,C,D,u\), the ordered triple (BP1),
and every predicate on these coordinates, and
\[
 c_F-c_H=2(A-1)\Lambda_H+6hj.
 \tag{BP4}
\]
All other full-range hits still have a computable image among
first-half hits after retaining the coordinate permutation. More
precisely, let \(\mathcal H_{\min}\) consist of first-half hit codes
whose first denominator is minimal and whose middle-channel lift has
\(y<z\). There is a surjection from full-range hits to
\(\mathcal H_{\min}\). For a fixed code in \(\mathcal H_{\min}\),
its entire inverse fibre is obtained by the following finite rule:
take all distinct permutations \((b,c,d)\) of its lifted triple,
retain those with
\[
 3h+1\le b\le9h,\qquad
 v_p((4b-p)c-pb)\in\{0,1\},
 \tag{BP5}
\]
and recover their full codes by the inverse in the proof.
Thus the two existential selectors agree, although their full hit
sets, shell labels and total counts need not agree.
\end{proposition}
\begin{proof}
For a full-range hit with \(a\le6h\), the conditions \(AB\mid a\)
give \(A,B\le6h\), and \(j<3h\). The hit tests on those same
coordinates are identical. The two proved radix decoders are inverse,
and the values \(\Lambda_F=3\Lambda_H\), \(18h-12h=6h\)
give (BP4).

For the second statement lift a full code to its ordered triple.
Choose its minimum denominator as first coordinate. Theorem
\ref{thm:bp-first-half} puts it in the first-half interval. The
residual-divisor proof in Theorem~\ref{thm:bp-code} gives a unique
canonical exterior orientation, with residual exponent 0 for the
second denominator, or a middle orientation with both exponents 1.
In the latter case choose the smaller of the final two denominators
as second. They cannot be equal: equality would give
\(d_y=d_z=S\), hence \(R\mid2S\), and the unit condition would
give \(R\mid2\), impossible. This yields a unique code in
\(\mathcal H_{\min}\). In the exterior case exactly one of the
two residual exponents is 0 and the other 2, so its orientation is
also unique. The minimum value and the multiset of remaining values
are unique, even if a value occurs more than once; hence these rules
give a unique resulting code.

For a retained permutation in (BP5), put \(R_b=4b-p,S_b=pb\) and
\(d_c=R_bc-S_b\). When \(v_p(d_c)=0\), recover
\(\varepsilon=0,u=b^2/d_c\); when \(v_p(d_c)=1\), recover
\(\varepsilon=1,u=pb^2/d_c\). The proof of Theorem
\ref{thm:bp-code} gives the positive integers \(A,B,C,D\) and
the code with \(j=b-3h-1\). Since \(A,B\le b\le9h\), it is
in the full rectangle and its ordered lift is exactly \((b,c,d)\).
Conversely every full code in the inverse fibre has that triple as
one of the distinct permutations, its first coordinate is in the
full interval, and its canonical residual exponent is 0 or 1.
Thus (BP5) gives every point of the fibre and no duplicates.
Surjectivity follows because each minimal first-half code has its
identical full-radix lift.
\end{proof}

\begin{corollary}[Boolean products and the exact two-layer inclusion]
Define
\[
 S(p)=1-\prod_{c=0}^{N_H-1}(1-\chi_{H,p}(c)),\qquad
 S_2(p)=1-\prod_{c=0}^{2\Lambda_H-1}(1-\chi_{H,p}(c)).
\]
Then \(S_2(p)\le S(p)\), \(S(p)=1\) if and only if
\(\operatorname{ES}(p)\), and \(S_2(p)=1\) if and only if a
first-half hit has \(A\in\{1,2\}\). For any prime set \(\mathcal P\)
define its infinite Boolean product as the infimum of its finite
initial products. Then
\[
 \prod_{p\in\mathcal P}S(p)=1
 \quad\Longleftrightarrow\quad
 (\forall p\in\mathcal P)\operatorname{ES}(p).
\]
\end{corollary}
\begin{proof}
Each complement factor is a bit. Its finite product is one exactly
when every hit bit vanishes. The radix intervals for fixed \(A\)
are \([(A-1)\Lambda_H,A\Lambda_H-1]\), so the first two are
exactly the displayed smaller range. Inclusion gives the inequality.
The finite initial products form a nonincreasing sequence of bits;
one zero factor makes every later product zero, and absent such a
factor every finite product is one. Apply (BP3).
\end{proof}

These statements correspond to transcript lines 14--163, 180--387,
5203--5293 and 6465--6606. The repository's previous exact radix map
is Theorem \texttt{raw-first-half-code} in
\nolinkurl{raw_scale_hit_incidence.tex}; the broader full-range
bound and code occur in \nolinkurl{bounded_transport.tex} and
\nolinkurl{unary_boolean_transport.tex}.

\subsection{Endpoint forcing and the corrected generalized carrier}

\begin{theorem}[The endpoint factor theorem]\label{thm:bp-endpoint}
Let \(p\equiv1\pmod{24}\) be prime. An \(A=1\) hit with
\(R^2\le p\) and \(B=1\) or \(D=1\) exists if and only if
\((p+1)/2\) or \(p+4\) has a prime divisor congruent to 3 modulo 4.
If \(q\) is the least such divisor of the first integer, the
exterior code is \((p-1)(q-3)/4\). If it is the least such divisor
of the second integer, the middle code is \((p+1)(q-1)/4\).
\end{theorem}
\begin{proof}
Both integers are odd and 1 modulo 4. If either has a prime factor
3 modulo 4, the number of such prime factors counted with multiplicity
is positive and even. The least one \(q\) therefore has square at
most that integer. For \((p+1)/2\), this gives \(q^2<p\).
For \(p+4\), it gives \(q^2\le p+4\); since \(q^2-p\)
is divisible by 8, it follows that \(q^2\le p\).

In the first case use
\[
 (R,A,B,D,\varepsilon)=
 (q,1,1,(p+q)/4,0),\quad C=(p+1)/q.
\]
In the second case use
\[
 (R,A,B,D,\varepsilon)=
 (q,1,(p+q)/4,1,1),\quad C=(p+q+4)/(4q).
\]
All these quantities are positive integers; in the second case oddness
of \(q\) and \(q\mid p+4\) give the last integrality. Substitution
in (BP2) gives the stated codes, and the tuple equations prove the hits.

Conversely \(A=B=1\) makes the middle channel impossible because
\(R\nmid2\); the exterior channel requires \(R\mid p+1\).
If \(A=D=1\), then \(B=a\), and the middle channel requires
\(R\mid p+4\). The exterior channel would require
\(R\mid p^2+4\). A positive integer \(R\equiv3\pmod4\)
has a prime factor \(q\equiv3\pmod4\). Such a prime cannot divide
\(p^2+4\): division by the unit 4 would make \(-1\) a square
modulo \(q\). A square root of \(-1\) would have order 4 in a
group of order \(q-1\equiv2\pmod4\), impossible. In each remaining
case that prime \(q\mid R\) divides \((p+1)/2\) or \(p+4\).
\end{proof}

The factor hypothesis is not universal: at \(p=193\), the two
integers are the primes 97 and 197, both 1 modulo 4. Nevertheless
\((A,B,C,D,\varepsilon)=(1,5,138,10,0)\) gives the code 200 and
the solution \((50,1380,1331700)\). This retains the actual scope of
transcript lines 2570--2944 and its correction at 5149--5189.

\begin{theorem}[Exact generalized carrier]
\label{thm:bp-crt}
For \(u>0\) set
\(\kappa(u)=\prod_{\ell^e\parallel u}\ell^{\lceil e/2\rceil}\).
Let \(n>0\) and \(R>0\), \(R\equiv3\pmod4\). The system
\[
 t\equiv1\pmod{12n},\quad
 t\equiv-R\pmod{4\kappa(u)},\quad
 t\equiv-R-4u\pmod{4R}
 \tag{BP6}
\]
has an integer solution if and only if
\[
 \gcd(3n,\kappa(u))\mid(R+1)/4,\qquad
 \gcd(3n,R)\mid(R+4u+1)/4.
 \tag{BP7}
\]
Its complete solution set, when nonempty, is one class modulo
\(4\operatorname{lcm}(3n,\kappa(u),R)\).
For the exterior carrier retain \(R\mid4u+1\) and omit the third
congruence; only the first divisibility in (BP7) remains.
For a prescribed prime \(p\), actual membership in the carrier,
\(R<p\), and the original common divisor must be retained.
\end{theorem}
\begin{proof}
Prime valuations give \(u\mid a^2\) if and only if
\(\kappa(u)\mid a\), and \(\kappa(u)\mid u\). With
\(4a=p+R\), these are exactly
\[
 u\mid a^2\ \Longleftrightarrow\ p\equiv-R\pmod{4\kappa(u)},
 \quad
 R\mid a+u\ \Longleftrightarrow\ p\equiv-R-4u\pmod{4R}.
\]
For a finite system of integer congruences, pairwise residue agreement
modulo the corresponding gcds is necessary by subtraction. To prove
sufficiency, for each prime select a modulus having maximal exponent
of that prime. Its prescribed residue agrees with every other at all
lower powers by the pairwise condition. Ordinary CRT over these distinct
prime powers gives a solution. Differences of solutions are exactly
multiples of the lcm. Applying this criterion to (BP6), the first two
pair conditions are (BP7); the third is
\(\gcd(\kappa(u),R)\mid u\), automatic because \(\kappa(u)\mid u\).
The lcm gives the stated period. The same argument with two moduli
gives the exterior statement.

For a class \(b+m\mathbb Z\), intersection with an integer interval
\([\alpha,\beta]\) occurs exactly when
\[
 \left\lceil\frac{\alpha-b}{m}\right\rceil
 \le\left\lfloor\frac{\beta-b}{m}\right\rfloor.
\]
At \(\alpha=\beta=p\), this is precisely membership of \(p\).
Finally a single \(u\mid a^2\) has the unique exponent vector
\((e_\ell)\in\prod_{\ell\mid a}\{0,\ldots,2v_\ell(a)\}\).
Reduction modulo different prime powers dividing \(R\) uses this same
vector. Its complete fibre is the intersection of the local fibres
on that box, not their Cartesian product.
\end{proof}

At \(n=2\), (BP7) is exactly
\[
 2\mid u\Rightarrow R\equiv7\pmod8,\quad
 3\mid u\Rightarrow R\equiv2\pmod3,\quad
 3\mid R\Rightarrow u\equiv2\pmod3.
\]
Thus a modulo-8 condition cannot be removed: \(R=3,u=2\) gives
incompatible requirements \(t\equiv1\pmod{24}\) and
\(t\equiv5\pmod8\). Also \(R=7,u=1\) gives the compatible
class \(73\pmod{168}\), which does not contain 193. These are the
corrections of transcript lines 4740--5058 and the generalization at
6104--6158.

\subsection{Cyclotomic order, infinitude, coverage and a smaller cofactor}

\begin{theorem}[Cyclotomic Euclid and complete-family coverage]
\label{thm:bp-cyclotomic}
Let \(m=12n\), \(n>0\), and \(P>0\). Then
\[
 \Phi_m(mP)>1,\qquad \gcd(mP,\Phi_m(mP))=1,
\]
and every prime divisor of this value is 1 modulo \(m\).
There are infinitely many primes in that class. For each such prime
\(p\), there are exactly \(\varphi(m)/2\) integers
\(1\le r\le(p-1)/2\) with \(p\mid\Phi_m(mr)\).
Consequently, with \(P_T=\prod_{r=1}^T\Phi_m(mr)\),
\[
 p\equiv1\pmod m,\quad p\le2T+1\quad\Longrightarrow\quad p\mid P_T.
 \tag{BP8}
\]
In the Boolean convention of the preceding corollary,
\[
 \prod_{p\equiv1\ (m)}S(p)
 =\inf_{T\ge1}\prod_{r=1}^T
   \prod_{\substack{q\mid\Phi_m(mr)\\q\ \text{prime}}}S(q).
 \tag{BP9}
\]
This identity does not assert that either side equals one.
\end{theorem}
\begin{proof}
Use \(X^m-1=\prod_{d\mid m}\Phi_d(X)\) and \(\Phi_m(0)=1\).
For a prime \(q\nmid mr\), a root \(r\) of \(\Phi_m\)
has order exactly \(m\): a smaller order \(e\mid m\) would make
it a common root with a factor \(\Phi_d\), \(d\mid e<m\),
and hence a repeated root of \(X^m-1\). This contradicts the
nonzero derivative \(mr^{m-1}\) modulo \(q\). Conversely an
element of order \(m\) vanishes on none of the proper-divisor
factors, so it vanishes on \(\Phi_m\).

For \(X=mP\), the constant term gives \(\Phi_m(X)\equiv1\pmod X\),
so each prime divisor is prime to \(mX\). Its order is \(m\),
which divides \(q-1\). Pair primitive roots into nonreal conjugates.
For real \(X\ge2\), every pair contributes \(|X-\zeta|^2\)
strictly between \((X-1)^2\) and \((X+1)^2\); hence
\[
 (X-1)^{\varphi(m)}<\Phi_m(X)<(X+1)^{\varphi(m)}.
\]
This proves \(\Phi_m(mP)>1\). If \(P\) is the product of any
finite list of primes 1 modulo \(m\), a prime divisor of that value
is outside the list and in the same class, proving infinitude.

For a prescribed prime \(p\equiv1\pmod m\), the cyclic group
\(\mathbb F_p^\times\) has exactly \(\varphi(m)\) elements
of order \(m\). Multiplication by the unit \(m^{-1}\) gives exactly
that many roots of \(\Phi_m(mX)\). As \(4\mid m\), negation
preserves primitive \(m\)-th roots: the exponent of a primitive
root changes from a unit \(k\pmod m\) to \(k+m/2\), still a
unit at every odd prime and still odd. Thus \(\Phi_m\) is even.
The nonzero representatives pair as \(r,p-r\), giving the exact
half-interval count and in fact
\[
 p^{\varphi(m)/2}\mid\prod_{r=1}^{(p-1)/2}\Phi_m(mr).
\]
This proves (BP8). Every prime occurring in the right side of (BP9)
is in the required class, and (BP8) eventually includes every member
of that class. Repetitions have no effect because \(S(q)^e=S(q)\)
for \(e\ge1\). Both infima therefore vanish for exactly the same
zero factors and otherwise equal one.
\end{proof}

The polynomial change needed for the larger modulus is exact:
\[
 \Phi_{12}(X^n)
 =\prod_{\substack{d\mid n\\\gcd(n/d,6)=1}}\Phi_{12d}(X),
 \quad
 \Phi_{12n}(X)=
 \prod_{\substack{d\mid n\\\gcd(d,6)=1}}
 (X^{4n/d}-X^{2n/d}+1)^{\mu(d)}.
 \tag{BP10}
\]
Indeed \(\operatorname{ord}(\zeta^n)=\operatorname{ord}(\zeta)/
\gcd(\operatorname{ord}(\zeta),n)=12\) precisely when
\(\operatorname{ord}(\zeta)=12d\), \(d\mid n\), and
\(\gcd(n/d,6)=1\). Simple roots give the first polynomial
identity; multiplicative M\"obius inversion gives the second.
Thus \(\Phi_{12n}(X)=X^{4n}-X^{2n}+1\) exactly when
\(n=2^\alpha3^\beta\). Any other prime divisor of \(n\)
gives an additional nonconstant factor in the first identity.
Changing only the argument does not impose order \(12n\):
\(\Phi_{12}(24)=13\cdot73\cdot349\).

The recursion \(P_0=1\), \(W_k=\Phi_m(mP_k)\),
\(P_{k+1}=P_kW_k\) has \(P_k=\prod_{i<k}W_i\) and pairwise
coprime \(W_i>1\), by the same constant-term calculation. It
therefore introduces a new prime at every step. Multiplying an input
also by a factorial \(b!\) excludes every prime at most \(b\).
These statements prove newness and, with such a factorial, strict
growth; (BP8)--(BP9) concern the entire family and do not assert
coverage by a particular single recursion. These are the mathematical
scopes of transcript lines 5438--6411 and 8031--8136.

\begin{theorem}[Short homogeneous \(\Phi_{12}\) factorization]
\label{thm:bp-short}
For every prime \(p\equiv1\pmod{12}\), there exist coprime
positive integers \(r,s\le\lfloor\sqrt p\rfloor\) such that
\[
 r^4-r^2s^2+s^4=ph,\qquad 1\le h<p.
 \tag{BP11}
\]
Every prime divisor of \(h\) is congruent to 1 modulo 12.
\end{theorem}
\begin{proof}
Choose \(w\) of order 12 modulo \(p\) and put
\(H=\lfloor\sqrt p\rfloor\). The map
\((i,j)\mapsto i-wj\) from \(\{0,\ldots,H\}^2\) into
\(\mathbb F_p\) has a repeated value because \((H+1)^2>p\).
The difference gives \((r_0,s_0)\ne(0,0)\),
\(|r_0|,|s_0|\le H\), \(r_0\equiv ws_0\pmod p\).
Neither coordinate vanishes: one zero would force the other to be
a multiple of \(p\) of absolute value less than \(p\).
Retain \(g=\gcd(|r_0|,|s_0|)\) and their two signs, and set
\(r=|r_0|/g,s=|s_0|/g\). These retained data recover the original
signed pair. Since \(g<p\), division by it is valid modulo \(p\).
Evenness in each coordinate gives
\(p\mid F(r,s):=r^4-r^2s^2+s^4\).

One has \(F=(r^2-s^2)^2+r^2s^2>0\). If \(r\ge s\),
then \(F=r^4-s^2(r^2-s^2)\le r^4\), and the opposite case is
symmetric. Thus \(0<F\le H^4<p^2\), giving (BP11).
If a prime \(q\mid F\) divided \(r\) or \(s\), it would
divide both, contrary to coprimality. A primitive pair gives
\(F\equiv1\pmod2\) and \(F\equiv1\pmod3\), so these two
primes are excluded. The unit \(rs^{-1}\) is a root of
\(\Phi_{12}\) modulo \(q\), and the order argument above gives
\(q\equiv1\pmod{12}\).
\end{proof}

This proves the assistant proposition at transcript lines 6612--6751
with the terminal alternative \(h=1\) retained. For comparison of
the cofactor and quadratic sign, there is an exact lift formula.
If \(F(r,s)=ph\), \(p\nmid rs\), then
\[
 \frac{F(r+pt,s)}p\equiv h+t\,2r(2r^2-s^2)\pmod p.
 \tag{BP12}
\]
The coefficient is a unit: \(2r^2=s^2\) would imply
\(F=3s^4/4=0\pmod p\), impossible. Hence unrestricted lifts
realize every quotient residue modulo \(p\), including either
nonzero square class. Primitivity can be retained by choosing \(t\)
to avoid the one forbidden class modulo each prime dividing \(s\),
using CRT together with the prescribed class modulo \(p\).
A positive sufficiently large representative gives a positive lift.
This exact statement does not retain the short-box bound.

\subsection{Necessary nonresidue factors and the exact character target}

\begin{lemma}[Divisor-of-a-sum obstruction]\label{lem:bp-square}
Let \(a,b,d>0\), \(\gcd(a,b)=1\), \(d\mid a+b\), and
\(d\equiv3\pmod4\). Then \(-d\) is not a square modulo \(4ab\).
\end{lemma}
\begin{proof}
Suppose \(t^2\equiv-d\pmod{4ab}\). The hypotheses give
\(\gcd(d,ab)=1\). For each odd prime \(\ell\mid ab\),
\(\left(\frac{-d}{\ell}\right)=1\). Quadratic reciprocity,
including its sign because \(d\equiv3\pmod4\), gives
\[
 \left(\frac\ell d\right)=\left(\frac{-d}\ell\right)=1.
\]
If \(2\mid ab\), reducing modulo 8 makes \(t\) odd and gives
\(d\equiv7\pmod8\), so \(\left(\frac2d\right)=1\).
Prime factorization and multiplicativity therefore imply
\(\left(\frac ad\right)=\left(\frac bd\right)=1\).
But \(a\equiv-b\pmod d\) gives
\(\left(\frac ad\right)=-\left(\frac bd\right)\), a contradiction.
\end{proof}

\begin{theorem}[A nonresidue prime in every normalized witness]
\label{thm:bp-nonresidue}
Every middle witness in (BP1) has a prime divisor \(\ell\mid AB\)
with \(\left(\frac\ell p\right)=-1\). Every exterior witness
has such a prime divisor of \(BC\).
\end{theorem}
\begin{proof}
For the middle channel, \(R\mid A+B\), \(\gcd(A,B)=1\),
and \(p\equiv-R\pmod{4AB}\). Lemma~\ref{lem:bp-square}
says \(p\) is not a square modulo \(4AB\). If every prime
\(\ell\mid AB\) were a residue modulo \(p\), reciprocity
would give \(\left(\frac p\ell\right)=1\) for each odd
such prime, because \(p\equiv1\pmod4\). A unit square root
modulo \(\ell\) lifts to every power: given \(x^2\equiv p\)
modulo \(\ell^e\), solve
\(2xt\equiv(p-x^2)/\ell^e\pmod\ell\) and replace
\(x\) by \(x+t\ell^e\). If \(2\mid AB\),
\(\left(\frac2p\right)=1\) and \(p\equiv1\pmod4\)
give \(p\equiv1\pmod8\). Every such unit is a square modulo
each 2-power: starting modulo 8, replacing an odd root \(x\)
modulo \(2^e\), \(e\ge3\), by \(x+t2^{e-1}\) toggles
the next square bit modulo \(2^{e+1}\). If \(AB\) is odd,
the needed modulo-4 root follows from \(p\equiv1\pmod4\).
CRT now makes \(p\) a square modulo \(4AB\), a contradiction.

For the exterior channel the tuple equation is
\[
 A(4BCD-1)=p(B+C).
\]
Since \(A<p\), set \(T=(B+C)/A\in\mathbb Z_{>0}\).
Then \(pT+1=4BCD\), \(T\equiv3\pmod4\), and \(T\mid B+C\).
A common prime divisor of \(B,C\) would divide \(A\), so
\(\gcd(B,C)=1\). Also \(p\nmid BC\): \(B<p\), and
\(p\mid C\) in the tuple equation would force \(p\mid A\).
If \(p\) were a square modulo \(4BC\), its inverse would be
one, contrary to Lemma~\ref{lem:bp-square} because
\(p^{-1}\equiv-T\pmod{4BC}\). The same prime-power and
reciprocity argument therefore supplies a nonresidue prime in \(BC\).
\end{proof}

For any fixed \(K\ge2\), put
\[
 m_K=\operatorname{lcm}\left(12n,8,
       \prod_{\substack{3\le\ell\le K\\\ell\ \text{prime}}}\ell\right).
\]
Every prime \(p\equiv1\pmod{m_K}\) makes all primes at most
\(K\) quadratic residues: use reciprocity for odd primes and the
modulo-8 supplementary law for 2. Theorem~\ref{thm:bp-nonresidue}
therefore forces a prime greater than \(K\) in the specified product
\(AB\) or \(BC\). In the middle channel, \(A\le2\) forces
that prime to divide \(B\), hence \(B>K\). Theorem
\ref{thm:bp-cyclotomic} produces infinitely many primes satisfying
these restrictions. This excludes fixed bounds on these pairs, not
all witnesses or the global two-layer assertion. These are the exact
quantifiers of transcript lines 6836--7265.

\begin{proposition}[A nonempty Jacobi-sign class and exact target counts]
\label{prop:bp-character}
Let \(p\ge13\) be prime and \(p\equiv1\pmod4\), and put
\(L=(p-1)/4\). There is a prime quadratic nonresidue
\(\ell\le L\). The integer
\(a=\ell\lceil(L+1)/\ell\rceil\) is a first-half shell denominator
containing a nonresidue prime factor. For this shell put
\[
 \tau=\prod_{r\mid a}(2v_r(a)+1),\qquad
 \sigma=\prod_{\substack{r\mid a\\(r/p)=1}}(2v_r(a)+1).
\]
Then the exact number of divisors \(u\mid a^2\) with
\(\left(\frac uR\right)=-1\) is
\((\tau-\sigma)/2\ge\tau/3>0\).

For every first-half shell, independently of this choice, let
\[
 W_a=\left\{\prod_{\ell\mid a}\ell^{e_\ell}\pmod R:
                   -v_\ell(a)\le e_\ell\le v_\ell(a)\right\}.
\]
Then \(W_a^{-1}=W_a\) and
\[
 E_a\cup M_a\ne\varnothing
 \quad\Longleftrightarrow\quad \{-1,-p\}\cap W_a\ne\varnothing.
 \tag{BP13}
\]
With \(\psi\) ranging over all characters of
\((\mathbb Z/R\mathbb Z)^\times\),
\[
 |E_a|+|M_a|=
 \frac1{\varphi(R)}\sum_\psi\psi(-1)(1+\psi(p))
       \prod_{\ell^v\parallel a}\sum_{e=-v}^v\psi(\ell)^e.
 \tag{BP14}
\]
\end{proposition}
\begin{proof}
If every integer 1 through \(L\) were a residue, their negatives
would also be residues because \(p\equiv1\pmod4\), exhausting
all \(2L\) nonzero residues. Every integer from \(L+1\) through
\(3L\) would then be a nonresidue. But \(2L=2\cdot L\) is a
product of two assumed residues and lies in that interval. The least
nonresidue is prime because otherwise all its prime factors are
smaller residues. The prescribed \(a\) satisfies
\(L+1\le a\le L+\ell\le2L\).

For an odd prime \(r\mid a\), reciprocity and
\(R\equiv-p\pmod r\), \(R\equiv3\pmod4\), give
\(\left(\frac rR\right)=\left(\frac rp\right)\): the two
reciprocity signs cancel. For \(r=2\), evenness of \(a\)
gives \(R\equiv-p\pmod8\) and the same identity by the
supplementary law. Summing the Jacobi symbols of divisors of \(a^2\)
therefore factors into geometric sums. A residue prime contributes
\(2v+1\), and a nonresidue prime contributes
\(\sum_{e=0}^{2v}(-1)^e=1\). The sum is \(\sigma\), the
total is \(\tau\), and the negative count follows. At least one
nonresidue factor gives \(\tau\ge3\sigma\).

Dividing \(u\mid a^2\) by \(a\) translates exponents by
\(-v_\ell(a)\), and inversion negates every centered exponent.
The middle target becomes \(-1\), while the exterior target
becomes \(-p^{-1}\); inversion changes it to \(-p\). This
proves (BP13). Orthogonality gives, for a unit \(t\),
\[
 \#\{u\mid a^2:u\equiv t\pmod R\}
 =\frac1{\varphi(R)}\sum_\psi\overline{\psi(t)}
       \prod_{\ell^v\parallel a}\sum_{e=0}^{2v}\psi(\ell)^e.
\]
Shift each exponent by \(-v\), retaining the factor \(\psi(a)\).
At \(t=-a\) its prefactor is \(\psi(-1)\); at
\(t=-4^{-1}\) it is \(\psi(-1)\psi(4a)=\psi(-1)\psi(p)\).
Adding gives (BP14). No individual summand is asserted nonnegative;
the total is an integer target count.
\end{proof}

There is a further exact comparison with the divisor-pair statistic
used in the dialogue. Let
\[
 T_{R,S}=\#\{(m,r):m,r\mid S,\ R\mid m+r\}.
\]
Character orthogonality gives
\[
 T_{R,S}=\frac1{\varphi(R)}\sum_\psi\psi(-1)
                  \left|\sum_{d\mid S}\psi(d)\right|^2.
 \tag{BP15}
\]
Indeed expansion has the target \(-mr^{-1}=1\) in the unit group.
Reduction by \(g=\gcd(m,r)\) maps these pairs onto the primitive
pairs \((m_0,r_0)\) with \(m_0,r_0\mid S\),
\(\gcd(m_0,r_0)=1\), \(R\mid m_0+r_0\). Its full fibre is
\[
 \{(g m_0,g r_0):g\mid S/(m_0r_0)\}.
 \tag{BP16}
\]
To prove this, division by \(g\mid S\) preserves the congruence
because \(g\) is a unit modulo \(R\). Conversely both scaled
entries divide \(S\) exactly when
\(g\mid\gcd(S/m_0,S/r_0)=S/(m_0r_0)\). The primitive pair
reconstructs the ordered solution
\[
 k=(m_0+r_0)/R,\qquad
 (a,y,z)=(a,(S/r_0)k,(S/m_0)k).
\]
Its reciprocal sum is \(1/a+R/S=4/p\). The inverse primitive
pair is the reduced positive ratio \((Ry-S)/S=y/z\); prime
valuations in \((Ry-S)(Rz-S)=S^2\) prove both its entries
divide \(S\). Thus \(T>0\) is exactly shell occupancy, while
its multiplicities are the divisor counts of (BP16). A canonical
hit supplies the distinct pair \((A,p^{1-\varepsilon}B)\),
so \(T\ge2\). This proves, with its full multiplicity map,
the relation in transcript lines 3594--3744 and 7843--7887.

\subsection{Five exact regressions and their scopes}

The following finite certificates are the five stronger-claim
counterexamples of transcript lines 8330--8786. They were reproduced
by a separate standard-library integer calculation. Their exhaustiveness
concerns the displayed finite divisor sets or coordinate boxes, not all
primes or all Erd\H{o}s--Straus candidates.

\paragraph{1. A selected nonresidue shell can be empty.}
Take \(p=3361\), \(L=840\). The primes 2,3,5,7 are residues
modulo \(p\): use \(p\equiv1\pmod8\), \(p\equiv1\pmod{12}\),
and reciprocity with \(p\equiv1\pmod5\), \(p\equiv1\pmod7\).
Thus every integer 1 through 10 is a residue. On the other hand
\(3361\equiv6\pmod{11}\), and the nonzero squares modulo 11
are \(\{1,3,4,5,9\}\), so 11 is the least nonresidue.
Proposition~\ref{prop:bp-character} selects
\[
 a=11\lceil841/11\rceil=847=7\cdot11^2,\quad R=27,\quad j=6.
\]
Every divisor of \(a^2\) is \(7^i11^e\), \(0\le i\le2\),
\(0\le e\le4\). Their residues are exactly
\[
 \{1,2,7,8,10,11,13,14,16,19,22,23,26\}\pmod{27}.
\]
They omit \(-4^{-1}=20\) and \(-a=17\), proving
\(E_6=M_6=\varnothing\). The six negative-Jacobi divisors are
\(11,77,539,1331,9317,65219\). The principal contribution
to (BP14) is \(2\cdot15/18=5/3\), while its complete sum is
zero; the other characters contribute \(-5/3\).
This is not an ES counterexample: the first shell has \(a=841\),
\(R=3\), \(u=29\mid841^2\), and \(3\mid4\cdot29+1\).
Its tuple is \((A,B,C,D,\varepsilon)=(29,1,1130,29,0)\), giving
\[
 c_H=79027200,\qquad
 \frac4{3361}=\frac1{841}+\frac1{950330}+\frac1{110139970}.
\]

\paragraph{2. Nonempty local fibres need not intersect in the divisor box.}
At \(p=37,a=18,R=35\), the complete divisor list is
\[
 \operatorname{Div}(324)=
 \{1,2,3,4,6,9,12,18,27,36,54,81,108,162,324\}.
\]
Those satisfying \(5\mid4u+1\) are \(1,6,36,81\); those
satisfying \(7\mid4u+1\) are \(12,54\). Both lists are nonempty
and their intersection is empty. CRT gives \(u\equiv26\pmod{35}\),
but no member of the original complete list belongs to that class.
This refutes the exchange of \(\forall m\exists u\) with
\(\exists u\forall m\), while preserving Theorem~\ref{thm:bp-crt}.

\paragraph{3. The short cofactor may always be one at a nonbase prime.}
The complete positive primitive boxes in Theorem~\ref{thm:bp-short}
give the following table:
\[
\begin{array}{c|c|c|c}
 p&H&\text{all primitive pairs}&h\\\hline
 61&7&(2,3),(3,2)&1\\
 73&8&(1,3),(3,1)&1\\
 193&13&(3,4),(4,3)&1\\
 2521&50&(1,26),(26,1)&181
\end{array}
\]
Here is a finite exact verification of completeness. The roots of
\(\Phi_{12}\) in these four prime fields are, respectively,
\[
 \{21,29,32,40\},\quad \{3,24,49,70\},\quad
 \{49,63,130,144\},\quad \{26,97,2424,2495\}.
\]
Substitution gives four distinct roots in each case, so the polynomial
degree proves the lists complete. For each root \(w\), the possible
positive coordinates in the box are exactly
\[
 r=ws-pk,\quad
 \max\left(1,\left\lceil\frac{pk+1}{w}\right\rceil\right)
 \le s\le
 \min\left(H,\left\lfloor\frac{pk+H}{w}\right\rfloor\right),
 \quad0\le k\le\lfloor wH/p\rfloor.
 \tag{BP17}
\]
Evaluating these integer intervals gives, before removing nonprimitive
pairs, respectively
\[
\begin{array}{c|l}
61&(2,3),(3,2),(4,6),(6,4)\\
73&(1,3),(2,6),(3,1),(6,2)\\
193&(3,4),(4,3),(6,8),(8,6),(9,12),(12,9)\\
2521&(1,26),(26,1).
\end{array}
\]
The gcd test and direct quartic evaluation give the preceding table.
For example \(2^4-2^2 3^2+3^4=61\). Hence \(h>1\) cannot
be demanded at every prime greater than 13 in this prescribed box.

\paragraph{4. A smaller cyclotomic prime need not be a nonresidue.}
The last row of that complete table gives
\[
 1-26^2+26^4=456301=2521\cdot181,
 \qquad 88^2=3\cdot2521+181.
\]
Thus \(181<2521\), \(181\equiv1\pmod{12}\), but
\(\left(\frac{181}{2521}\right)=1\). Both permitted pairs
give that cofactor. The actual middle tuple
\((A,B,C,D,\varepsilon)=(2,159,7,2,1)\) has
\(AB=2\cdot3\cdot53\), with 53 a nonresidue modulo 2521:
reciprocity gives
\[
 \left(\frac{53}{2521}\right)
 =\left(\frac{30}{53}\right)
 =\left(\frac2{53}\right)\left(\frac3{53}\right)
  \left(\frac5{53}\right)=(-1)(-1)(-1)=-1.
\]
Its code and original denominators are
\[
 c_H=1600517,\qquad
 \frac4{2521}=\frac1{636}+\frac1{70588}+\frac1{5611746}.
\]
Substitution gives \(4ABCD=17808=A+B+pC\), and all the code
bits equal one. Thus the cofactor and witness factor have their
respective exact roles.

\paragraph{5. The degree-four cofactor bound does not persist for
\(\Phi_{24}\).}
For \(p=73\), the eight roots of \(\Phi_{24}(X)=X^8-X^4+1\)
modulo 73 are \(7,17,21,30,43,52,56,66\). Substitution and degree
prove completeness. Formula (BP17), now with \(H=8\) and this root
list, gives exactly
\((1,7),(7,1),(4,5),(5,4)\). They are all primitive, and
\[
 F_{24}(1,7)=5762401=73\cdot78937,\quad
 F_{24}(4,5)=296161=73\cdot4057.
\]
The smallest quotient is \(4057>73\), disproving the unchanged
requirement \(h<p\). The valid bound for this degree-eight trinomial
is \(F_{24}(r,s)\le\max(r,s)^8<p^4\), hence \(h<p^3\).
The cyclotomic infinitude theorem at modulus 24 is unaffected.

The primality checks needed in these finite certificates are elementary
trial divisions by all primes at most the respective square roots.
For 61 and 73 those primes are \(2,3,5,7\); for 193 they are
\(2,3,5,7,11,13\); for 2521 they are
\(2,3,5,7,11,13,17,19,23,29,31,37,41,43,47\); for 3361
add 53 to that last list. None divides the corresponding number.
For the auxiliary primes 97,181,197 the same test uses respectively
the lists through 7,13,13. The complete finite divisor sets, root
lists, quotient values and reciprocal identities above are exact
certificates for their stated counterexamples. They do not establish
the universal positivity of (BP3) or (BP9).
